% dir-callers.tex -- 移行済みクラス呼出しの総合ゲート。
%   * 脚注合印（\@makefnmark d 分岐: \rlap + 箱工場）
%   * maketitle 著者表組み（\if@tate 時の \ialign への \tate shim は
%     d では適用しない: 内枠は既に正しい向き）
%   * theindex（d: ページ方向フラグ + \hsize=\textheight の役割交換で
%     本当の横組みページを出し、終端で tate 復帰）
%   * complements の theindex_yoko 既定（tate 文書では false = tate 索引）
\documentclass[pdflatex,tate,paper=a5]{jlreq}
\makeatletter
\begin{document}
% --- CALLER-FOOTNOTE: 脚注（合印は横組で、本文は縦組）
本文\footnote{脚注は縦組み}の続き。
\typeout{TATE-PASS: CALLER-FOOTNOTE (footnote mark compiled via factory)}%
% --- CALLER-MAKETITLE: 著者表組み（tabular の行は右から左の列になる）
\title{縦組標題}
\author{著者一 \\ 著者二}
\date{}
\maketitle
\typeout{TATE-PASS: CALLER-MAKETITLE (author tabular compiled under d)}%
\newpage
% --- CALLER-INDEX: 索引（jlreq.cls の theindex: 常に横書きページ）
本文。
\begin{theindex}
\item 阿伊, 3
\end{theindex}
% --- CALLER-INDEX-RESTORE: 索引後に縦組みに戻っている
本文再開。
\ifnum\ptj@dir=1
  \typeout{TATE-PASS: CALLER-INDEX-RESTORE (dir=1 after theindex)}%
\else
  \typeout{TATE-FAIL: CALLER-INDEX-RESTORE (dir=\the\ptj@dir, expected 1)}%
\fi
% complements 既定の確認: tate 文書では theindex_yoko=false（索引自体は
% tate のまま。uplatex 同一挙動）。横組み索引ページは theindex_yoko=true
% 時の dir-index-yoko fixture が検査する。
\end{document}
